%%% =======================================================================
%%% ISLS Extension Phase 10 v1.0.0 (Revised v2)
%%% Part A: C28 Babylon Bridge — IR-Validated Multi-Language Forge
%%% Part B: Benchmark Update — B16–B24 for Forge/Foundry/Oracle/Studio
%%% Author: Sebastian Klemm
%%% Date: 18 March 2026
%%% Status: Implementation Specification for Claude Code
%%% =======================================================================
\documentclass[11pt,a4paper,twoside]{article}

\usepackage[utf8]{inputenc}
\usepackage{amsmath,amssymb,amsthm}
\usepackage{booktabs,longtable,tabularx}
\usepackage{enumitem}
\usepackage{geometry}
\usepackage{hyperref}
\usepackage{cleveref}
\usepackage{xcolor}
\usepackage{fancyhdr}
\usepackage{listings}
\usepackage{titlesec}

\geometry{margin=2.5cm,top=3cm,bottom=3cm,headheight=14pt}

\theoremstyle{definition}
\newtheorem{definition}{Definition}[section]
\newtheorem{requirement}{Requirement}[section]

\theoremstyle{plain}
\newtheorem{theorem}{Theorem}[section]

\theoremstyle{remark}
\newtheorem{remark}{Remark}[section]

\newcommand{\ISLS}{\textsf{ISLS}}

\definecolor{sec-color}{HTML}{1e3a5f}

\titleformat{\section}
  {\Large\bfseries\color{sec-color}}
  {\thesection}{1em}{}

\lstset{
  basicstyle=\ttfamily\small,
  breaklines=true,
  frame=single,
  backgroundcolor=\color{gray!5},
  rulecolor=\color{gray!30},
}

\pagestyle{fancy}
\fancyhf{}
\fancyhead[LE,RO]{\ISLS{} Phase 10 v1.0.0 (Rev.\ 2)}
\fancyhead[RE,LO]{Sebastian Klemm}
\fancyfoot[C]{\thepage}

\hypersetup{
  colorlinks=true,
  linkcolor=blue!60!black,
  pdftitle={ISLS Phase 10 Rev 2 — Babylon Bridge + Benchmarks},
  pdfauthor={Sebastian Klemm},
}

\begin{document}

\begin{titlepage}
\centering
\vspace*{2cm}
{\Huge\bfseries \ISLS{} Extension Phase~10\\[0.3cm]
v1.0.0 (Revised v2)}\\[1.5cm]
{\Large Part A: C28 --- Babylon Bridge\\
Part B: Benchmark Update B16--B24}\\[1cm]
{\large Honest about what exists.\\
Honest about what must be built.\\
Honest about what must be measured.}\\[2cm]
{\large Sebastian Klemm}\\[0.5cm]
{\large 18 March 2026}

\begin{abstract}
\noindent
\textbf{Part A} integrates the Babylon Compiler's IR layer
(\texttt{glyph-ir}, \texttt{glyph-embed}, \texttt{glyph-canon})
as a structural validation and type-consistency backbone for the
Forge. The Oracle (C25) still generates code --- but now against
an IR-derived structural contract. Codegen backends are \textbf{not}
claimed to exist in Babylon or Merkaba; they are specified here as
new code that must be written. Six backends (Rust, TypeScript,
Python, SQL, YAML/Docker, Markdown) emit from IR structure + Oracle
content.

\textbf{Part B} extends the benchmark suite from B01--B15 (core
pipeline) to B16--B24, covering: Oracle synthesis latency, Oracle
rejection rate, Pattern Memory hit rate, Foundry compilation success
rate, Foundry average attempts, Template match accuracy, Studio
endpoint latency, full forge-to-project throughput, and autonomy
ratio over time.

\medskip\noindent\textbf{Status:} Implementation Specification.\\
\textbf{New Crate:} C28 (\texttt{isls-multilang}).\\
\textbf{Depends on:} C22, C23, C25, C26, C27;
  Babylon: \texttt{glyph-ir}, \texttt{glyph-embed},
  \texttt{glyph-canon}, \texttt{glyph-q16}.
\end{abstract}
\end{titlepage}

\tableofcontents
\newpage


%%% =====================================================================
%%% PART A
%%% =====================================================================
\part{Part A: C28 --- Babylon Bridge}\label{part:bridge}

\section{What Babylon Actually Provides}\label{sec:provides}

An honest inventory of what exists in the Babylon Compiler repository
and what does \textbf{not} exist:

\begin{tabularx}{\textwidth}{l c X}
\toprule
\textbf{Component} & \textbf{Exists?} & \textbf{Detail} \\
\midrule
\texttt{glyph-ir} & \checkmark & IrDocument with 10 NodeKinds, 11
  EdgeKinds, canonical sort, SHA-256 digest. \\
\texttt{glyph-embed} & \checkmark & 5-axis H5 embedding: structural
  coupling, functional density, topological complexity, symmetry,
  entropy. Q16 fixed-point. \\
\texttt{glyph-canon} & \checkmark & JCS (RFC~8785) canonical JSON +
  SHA-256. \\
\texttt{glyph-q16} & \checkmark & Q16.16 fixed-point arithmetic. \\
\texttt{glyph-gate} & \checkmark & Obligation evaluation, tripolar
  gate, PoR FSM. \\
\texttt{glyph-mef} & \checkmark & Append-only hash chain for
  compilation evidence. \\
\texttt{glyph-frontends} & \checkmark & 3 parsers (Sanskroot, HanLan,
  Cuneiform) lowering to glyph-ir. \\
\midrule
Rust codegen backend & \texttimes & Does not exist. \\
TypeScript codegen backend & \texttimes & Does not exist. \\
Python codegen backend & \texttimes & Does not exist. \\
SQL codegen backend & \texttimes & Does not exist. \\
Any codegen backend & \texttimes & Babylon emits Crystal bundles,
  not source code in target languages. \\
Merkaba codegen backends & \texttimes & Mentioned in concepts but
  not implemented as functional emitters. \\
\bottomrule
\end{tabularx}

\begin{remark}
Babylon is a compiler that \emph{reads} source code and produces
certified knowledge artifacts. It does not \emph{write} source code.
The codegen backends must be written from scratch as part of C28.
\end{remark}


\section{Integration Architecture}\label{sec:arch}

\begin{lstlisting}[title={The Honest Pipeline}]
DecisionSpec
  |
  v
[PMHD Drill] -> PmhdMonolith             (existing, C21)
  |
  v
[ArtifactIR] -> Components + Interfaces   (existing, C22)
  |
  v
[Bridge: ArtifactIR -> glyph-ir]          (NEW, C28)
  |   Converts components to IrNodes (Module, Function)
  |   Converts interfaces to IrEdges (CalleeRef, Contains)
  |   Deterministic, content-addressed
  v
[glyph-embed: compute H5 embedding]       (existing, Babylon)
  |   5 axes validate structural coherence
  |   Catches: disconnected modules, over-coupling,
  |   fragmentation, asymmetry, repetition
  v
[Oracle fills function bodies]             (existing, C25)
  |   Oracle receives: function name, signature, interfaces,
  |   constraints, H5 context
  |   Oracle returns: implementation code STRING for each function
  |   The Oracle does NOT return IR. It returns code.
  v
[Codegen Backend assembles files]          (NEW, C28)
  |   Takes: IR structure + Oracle body strings
  |   Produces: complete source files with correct
  |   imports, module structure, type definitions
  |   The backend writes the SCAFFOLDING.
  |   The Oracle writes the LOGIC.
  v
[Foundry: compile + test + fix loop]       (existing, C27)
  |
  v
[Crystal with IR digest + H5 embedding]
\end{lstlisting}

\begin{remark}[Key Honesty]
The Oracle still generates code strings. The codegen backends do
\textbf{not} generate logic --- they generate scaffolding (imports,
module declarations, type signatures, function stubs, test
scaffolding). The Oracle fills the function bodies. The backend
assembles the final files from scaffolding + Oracle content.

This is a meaningful improvement over pure Oracle generation because:
(a)~imports are always correct (derived from IR edges),
(b)~type signatures are always consistent (derived from IR nodes),
(c)~module structure is always coherent (derived from IR topology),
(d)~only the \emph{bodies} come from the Oracle, not the entire file.
\end{remark}


\section{ArtifactIR to glyph-ir Conversion}\label{sec:convert}

\begin{definition}[Bridge Conversion]\label{def:bridge}
\begin{lstlisting}
pub fn artifact_to_glyph_ir(ir: &ArtifactIR) -> glyph_ir::IrDocument {
    let mut doc = glyph_ir::IrDocument::new(
        &ir.header.domain,
        &hex_encode(&ir.header.artifact_id),
    );

    // Root module
    let root = glyph_ir::IrNode::new("n_root", NodeKind::Module, "root");
    doc.nodes.push(root);

    // Each component -> Function node with metadata
    for (i, comp) in ir.components.iter().enumerate() {
        let id = format!("n_fn_{}", i);
        let mut node = glyph_ir::IrNode::new(&id, NodeKind::Function, &comp.name);
        node.params = if comp.dependencies.is_empty() {
            None
        } else {
            Some(comp.dependencies.clone())
        };
        // Store component content as property for Oracle prompt
        let mut props = BTreeMap::new();
        props.insert("kind".into(), json!(comp.kind));
        props.insert("description".into(), json!(comp.content));
        node.properties = Some(props);
        doc.nodes.push(node);
        doc.edges.push(glyph_ir::IrEdge::new("n_root", &id, EdgeKind::Contains));
    }

    // Interfaces -> edges
    for iface in &ir.interfaces {
        let provider_id = find_node_by_name(&doc, &iface.provider);
        let consumer_id = find_node_by_name(&doc, &iface.consumer);
        if let (Some(p), Some(c)) = (provider_id, consumer_id) {
            doc.edges.push(glyph_ir::IrEdge::new(&p, &c, EdgeKind::CalleeRef));
        }
    }

    doc.canonicalize();
    doc
}
\end{lstlisting}
\end{definition}


\section{Codegen Backends (Must Be Built)}\label{sec:codegen}

\begin{definition}[CodegenBackend Trait]\label{def:backendtrait}
\begin{lstlisting}
pub trait CodegenBackend: Send + Sync {
    fn language(&self) -> &str;
    fn extension(&self) -> &str;

    /// Generate scaffolding from IR structure.
    /// Returns files with stubs where Oracle content will be inserted.
    fn scaffold(&self, doc: &glyph_ir::IrDocument) -> Vec<ScaffoldFile>;

    /// Assemble final files by merging scaffolding with Oracle bodies.
    fn assemble(
        &self,
        scaffolds: &[ScaffoldFile],
        oracle_bodies: &BTreeMap<String, String>,  // node_id -> code
    ) -> Vec<EmittedFile>;
}

pub struct ScaffoldFile {
    pub path: String,
    pub language: String,
    pub imports: Vec<String>,         // derived from IR edges
    pub type_definitions: Vec<String>, // derived from IR nodes
    pub function_stubs: Vec<FunctionStub>,
}

pub struct FunctionStub {
    pub node_id: String,       // glyph-ir node ID
    pub name: String,
    pub signature: String,     // generated by backend from IR
    pub body_placeholder: String, // "// TODO: Oracle fills this"
}

pub struct EmittedFile {
    pub path: String,
    pub content: String,
    pub language: String,
    pub scaffold_lines: usize,  // lines from backend (structure)
    pub oracle_lines: usize,    // lines from Oracle (logic)
}
\end{lstlisting}
\end{definition}

\begin{definition}[Six Backends]\label{def:backends}
\begin{tabularx}{\textwidth}{l l l X}
\toprule
\textbf{Backend} & \textbf{Lang} & \textbf{Ext} & \textbf{Scaffolding It Generates} \\
\midrule
\texttt{RustBackend} & Rust & .rs & \texttt{mod} declarations,
  \texttt{pub struct} with fields, \texttt{pub fn} signatures,
  \texttt{use} imports from IR edges, Cargo.toml. \\
\texttt{TypeScriptBackend} & TS & .ts & \texttt{export interface} for
  types, \texttt{export function} signatures, \texttt{import} from IR
  edges, package.json. \\
\texttt{PythonBackend} & Python & .py & \texttt{@dataclass class} for
  types, \texttt{def} signatures with type hints, \texttt{import} from
  IR edges, requirements.txt. \\
\texttt{SqlBackend} & SQL & .sql & \texttt{CREATE TABLE} from type
  nodes, \texttt{CREATE INDEX} from edges, migration numbering. \\
\texttt{OpsBackend} & YAML & .yml & Dockerfile skeleton,
  docker-compose services, CI pipeline steps. Structure-only,
  no Oracle needed. \\
\texttt{DocsBackend} & Markdown & .md & README from IR root node,
  module docs from Function nodes, API docs from interface edges.
  Structure-only. \\
\bottomrule
\end{tabularx}

\textbf{Key distinction:} OpsBackend and DocsBackend are
\textbf{fully deterministic} --- they produce complete files from the
IR alone, no Oracle needed. RustBackend, TypeScriptBackend,
PythonBackend, SqlBackend produce scaffolding that needs Oracle
bodies.
\end{definition}


\section{Cross-Language Type Consistency}\label{sec:crosslang}

\begin{definition}[Type Node Convention]\label{def:typenode}
A type is represented as an IR Function node with
\texttt{kind = ``type''} in its properties and fields listed as
parameters:
\begin{lstlisting}
// IR representation of a "User" type:
IrNode {
    id: "n_type_user",
    kind: Function,  // glyph-ir's Function serves as container
    name: "User",
    params: Some(vec!["id:i64", "name:String", "email:String"]),
    properties: Some({"kind": "type", "description": "User entity"}),
}
\end{lstlisting}

Each backend interprets this IR node according to its language:
\begin{lstlisting}
// RustBackend:
pub struct User { pub id: i64, pub name: String, pub email: String }

// TypeScriptBackend:
export interface User { id: number; name: string; email: string; }

// PythonBackend:
@dataclass
class User:
    id: int
    name: str
    email: str

// SqlBackend:
CREATE TABLE users (
    id BIGINT PRIMARY KEY,
    name TEXT NOT NULL,
    email TEXT NOT NULL
);
\end{lstlisting}

Consistency is structural: all backends read the same IR node. A
field name mismatch is impossible because there is only one source.
\end{definition}


\section{Embedding as Pre-Codegen Validation}\label{sec:h5}

\begin{definition}[Structural Pre-Check]\label{def:precheck}
Before codegen, the H5 embedding is computed:
\begin{lstlisting}
let embedding = glyph_embed::compute_embedding(&ir_doc);
\end{lstlisting}

Range checks per axis (configurable, defaults):
\begin{itemize}[nosep]
  \item $a_1$ (structural coupling): $[0.05, 0.80]$. Zero = disconnected. High = spaghetti.
  \item $a_2$ (functional density): $[0.10, 0.90]$. Zero = no functions. High = fragmented.
  \item $a_3$ (topological complexity): $[0.0, 0.70]$. High = over-nested.
  \item $a_4$ (symmetry): $[0.20, 1.0]$. Low = lopsided modules.
  \item $a_5$ (entropy): $[0.10, 0.90]$. Low = repetitive. High = incoherent.
\end{itemize}

If any axis is outside range, the system logs a warning and
optionally re-decomposes (retry with different PMHD strategy).
This is not a hard gate --- it is diagnostic. The Foundry's
compile/test loop is the hard gate.
\end{definition}


\section{Full-Stack Templates}\label{sec:templates}

Six new templates added to C26's catalog. Each is an IR tree (not a
prompt template) with multi-language atoms:

\begin{tabularx}{\textwidth}{l l r r X}
\toprule
\textbf{ID} & \textbf{Name} & \textbf{Atoms} & \textbf{Mol.} &
  \textbf{Languages} \\
\midrule
T11 & SaaS Starter & 22 & 6 & Rust + TS + SQL + Docker + MD \\
T12 & Dashboard & 18 & 5 & Rust + TS + SQL + Docker \\
T13 & API + Docs & 12 & 3 & Rust + MD + Docker \\
T14 & Python ML & 10 & 3 & Python + SQL + Docker \\
T15 & Static Site & 8 & 2 & HTML/CSS + MD \\
T16 & Monorepo & 16 & 5 & Rust + TS + SQL + Docker + MD \\
\bottomrule
\end{tabularx}

Each template atom specifies its backend:
\begin{lstlisting}
pub struct MultiLangAtom {
    pub name: String,
    pub backend: String,       // "rust", "typescript", "python", etc.
    pub fill: FillStrategy,    // Oracle, Static, Derive
    pub ir_node_ids: Vec<String>, // which IR nodes this atom covers
}
\end{lstlisting}


\section{CLI Extension}\label{sec:cli}

\begin{lstlisting}
# Forge with explicit target language
isls forge --spec intent.json --lang rust
isls forge --spec intent.json --lang typescript
isls forge --spec intent.json --lang python

# Forge full-stack (template auto-selects languages)
isls forge --spec intent.json --template saas-starter

# Emit IR for inspection
isls forge --spec intent.json --lang rust --dump-ir output.json

# Validate IR structure via H5 embedding
isls babylon check --ir output.json
\end{lstlisting}


\section{Acceptance Tests (Part A)}\label{sec:testacc}

\begin{longtable}{l l p{5.5cm}}
\toprule
\textbf{ID} & \textbf{Name} & \textbf{Procedure} \\
\midrule
\endfirsthead \midrule \endhead \bottomrule \endfoot
AT-BB1 & IR conversion & Convert ArtifactIR to IrDocument; verify
  node count = component count + 1 (root). \\
AT-BB2 & IR determinism & Convert same ArtifactIR twice; verify
  identical digest. \\
AT-BB3 & H5 embedding & Compute embedding; verify all 5 axes in
  Q16 range. \\
AT-BB4 & Rust scaffold & Generate scaffold from IR; verify
  .rs files with correct \texttt{pub fn} signatures. \\
AT-BB5 & TS scaffold & Generate scaffold; verify .ts files with
  correct \texttt{export interface} definitions. \\
AT-BB6 & Python scaffold & Generate scaffold; verify .py files
  with correct \texttt{@dataclass} definitions. \\
AT-BB7 & SQL scaffold & Generate scaffold; verify .sql files
  with correct \texttt{CREATE TABLE}. \\
AT-BB8 & Cross-lang types & Define type in IR; scaffold to Rust +
  TS; verify field names match exactly. \\
AT-BB9 & Assembly & Scaffold + mock Oracle bodies; assemble;
  verify complete files with both scaffold and body content. \\
AT-BB10 & Ops backend & Generate Dockerfile + compose from IR;
  verify valid YAML (no Oracle needed). \\
AT-BB11 & Docs backend & Generate README from IR; verify
  contains module names and descriptions. \\
AT-BB12 & Full-stack T11 & Forge SaaS Starter; verify dirs for
  all 5 languages exist with correct files. \\
\end{longtable}


%%% =====================================================================
%%% PART B
%%% =====================================================================
\part{Part B: Benchmark Update B16--B24}\label{part:bench}

\section{Current State}\label{sec:benchstate}

The benchmark suite (C10 \texttt{isls-harness}) currently covers
B01--B15:

\begin{tabularx}{\textwidth}{l l r l}
\toprule
\textbf{ID} & \textbf{Name} & \textbf{Last Value} & \textbf{Phase} \\
\midrule
B01 & Ingestion throughput & 234,897 obs/sec & Core \\
B02 & Graph update scaling & 156 $\mu$s & Core \\
B03 & Extraction scaling & 103 ms & Core \\
B04 & Cascade time & --- & Core \\
B05 & Dual consensus overhead & 133\% & Core \\
B06 & Crystal serialization & 55 $\mu$s & Core \\
B07 & Replay speed & 18,111 steps/sec & Core \\
B08 & Evidence verification & 5.0 $\mu$s & Core \\
B09 & Memory scaling (5K ent.) & 0.99 MB & Core \\
B10 & Full macro-step & --- & Core \\
B11 & Registry resolution & 0.07 $\mu$s & Phase~1 \\
B12 & Manifest construction & 4 ms & Phase~1 \\
B13 & Manifest verification & 3 ms & Phase~1 \\
B14 & Capsule roundtrip & 10.1 $\mu$s & Phase~1 \\
B15 & Scheduler overhead & 0\% & Phase~1 \\
\bottomrule
\end{tabularx}

\textbf{Missing:} Phases~5--9 (Forge, Oracle, Templates, Foundry,
Studio) have zero benchmarks. The harness has no metrics for
generative pipeline performance.


\section{New Benchmarks B16--B24}\label{sec:newbench}

\begin{longtable}{l l l p{5cm}}
\toprule
\textbf{ID} & \textbf{Name} & \textbf{Unit} & \textbf{What It Measures} \\
\midrule
\endfirsthead \midrule \endhead \bottomrule \endfoot

B16 & \texttt{forge\_throughput} & crystals/sec & End-to-end: DecisionSpec
  $\to$ Crystal, averaged over 10 forge runs with test specs. Includes
  PMHD drill + ArtifactIR + Oracle mock + Forge pipeline. \\

B17 & \texttt{oracle\_latency} & ms/call & Average Oracle synthesis
  call latency. With mock Oracle (measures overhead, not API latency)
  and optionally with real API (measures end-to-end). \\

B18 & \texttt{oracle\_rejection\_rate} & percent & Fraction of Oracle
  outputs that fail validation (parse, constraint, PMHD, gate). With
  mock Oracle: should be 0\%. Tracked as M34. \\

B19 & \texttt{pattern\_hit\_rate} & percent & Fraction of synthesis
  requests served from Pattern Memory without Oracle call. Initially
  0\%; grows over time. Tracked as M33 (autonomy ratio). \\

B20 & \texttt{foundry\_compile\_rate} & percent & Fraction of Foundry
  fabrication runs where first attempt compiles. Measures Oracle
  code quality. In mock mode: 100\% (mock always passes). In real
  mode: tracked per language. \\

B21 & \texttt{foundry\_avg\_attempts} & attempts & Average number of
  compile-fix cycles per successful fabrication. Target: $\le 2$.
  Measures Oracle+Foundry loop efficiency. \\

B22 & \texttt{template\_match\_accuracy} & percent & Fraction of
  test intents where auto-template-match selects the correct template.
  Measured against a labeled test set of 20 intents. \\

B23 & \texttt{gateway\_latency} & ms/request & Average round-trip
  for key Gateway endpoints: \texttt{GET /health},
  \texttt{GET /crystals}, \texttt{POST /forge}. \\

B24 & \texttt{full\_fabrication\_time} & seconds & Wall-clock time
  from DecisionSpec to a project on disk that passes
  \texttt{cargo test}. The ultimate end-to-end metric. \\

\end{longtable}


\section{Benchmark Implementation}\label{sec:benchimpl}

\begin{lstlisting}[title={New benchmark functions in isls-harness}]
pub fn bench_forge_throughput(engine: &mut ForgeEngine) -> BenchmarkResult {
    let specs = generate_test_specs(10); // 10 diverse test specs
    let start = Instant::now();
    let mut crystals = 0;
    for spec in &specs {
        if engine.forge(spec.clone()).is_ok() {
            crystals += 1;
        }
    }
    let elapsed = start.elapsed();
    BenchmarkResult {
        id: "B16".into(),
        name: "forge_throughput".into(),
        value: crystals as f64 / elapsed.as_secs_f64(),
        unit: "crystals/sec".into(),
    }
}

pub fn bench_oracle_latency(oracle: &mut OracleEngine) -> BenchmarkResult {
    let test_prompts = generate_test_prompts(20);
    let start = Instant::now();
    for prompt in &test_prompts {
        let _ = oracle.synthesize_raw(prompt);
    }
    let elapsed = start.elapsed();
    BenchmarkResult {
        id: "B17".into(),
        name: "oracle_latency".into(),
        value: elapsed.as_millis() as f64 / 20.0,
        unit: "ms/call".into(),
    }
}

pub fn bench_template_match(catalog: &TemplateCatalog) -> BenchmarkResult {
    let test_set = labeled_intent_test_set(); // 20 (intent, expected_template) pairs
    let mut correct = 0;
    for (intent, expected) in &test_set {
        let matched = catalog.best_match_for_intent(intent);
        if matched.as_ref().map(|t| &t.name) == Some(expected) {
            correct += 1;
        }
    }
    BenchmarkResult {
        id: "B22".into(),
        name: "template_match_accuracy".into(),
        value: correct as f64 / test_set.len() as f64 * 100.0,
        unit: "percent".into(),
    }
}

pub fn bench_full_fabrication(foundry: &mut Foundry) -> BenchmarkResult {
    let spec = simple_rust_api_spec(); // "health check endpoint"
    let start = Instant::now();
    let result = foundry.fabricate(FoundryOperation::NewProject {
        spec,
        output_dir: temp_dir(),
    });
    let elapsed = start.elapsed();
    BenchmarkResult {
        id: "B24".into(),
        name: "full_fabrication_time".into(),
        value: elapsed.as_secs_f64(),
        unit: "seconds".into(),
    }
}
\end{lstlisting}


\section{Report Integration}\label{sec:report}

The HTML report gains a new section: \textbf{``11. Generative Pipeline
Benchmarks (B16--B24)''} with the same layout as the existing
benchmark table. The Studio Dashboard gains a ``Forge Performance''
panel showing B16, B21, B24 as headline numbers.

\begin{lstlisting}[title={Report Section 11}]
<h2>11. Generative Pipeline Benchmarks</h2>
<table>
  <tr><td>B16</td><td>forge_throughput</td>
      <td>{value}</td><td>crystals/sec</td></tr>
  <tr><td>B17</td><td>oracle_latency</td>
      <td>{value}</td><td>ms/call</td></tr>
  <tr><td>B18</td><td>oracle_rejection_rate</td>
      <td>{value}</td><td>percent</td></tr>
  <tr><td>B19</td><td>pattern_hit_rate</td>
      <td>{value}</td><td>percent</td></tr>
  <tr><td>B20</td><td>foundry_compile_rate</td>
      <td>{value}</td><td>percent</td></tr>
  <tr><td>B21</td><td>foundry_avg_attempts</td>
      <td>{value}</td><td>attempts</td></tr>
  <tr><td>B22</td><td>template_match_accuracy</td>
      <td>{value}</td><td>percent</td></tr>
  <tr><td>B23</td><td>gateway_latency</td>
      <td>{value}</td><td>ms/request</td></tr>
  <tr><td>B24</td><td>full_fabrication_time</td>
      <td>{value}</td><td>seconds</td></tr>
</table>
\end{lstlisting}


\section{Benchmark CLI}\label{sec:benchcli}

\begin{lstlisting}
# Run all benchmarks (B01-B24)
isls bench

# Run only generative benchmarks (B16-B24)
isls bench --suite generative

# Run only core benchmarks (B01-B15)
isls bench --suite core

# Run a specific benchmark
isls bench --id B24

# Run with real Oracle (API key required)
isls bench --suite generative --oracle live
\end{lstlisting}


%%% =====================================================================
\part{Handoff}\label{part:handoff}

\section{Dependency Map}

\begin{lstlisting}
isls-multilang (C28)
  depends on:
    isls-types, isls-artifact-ir (C22), isls-forge (C23),
    isls-oracle (C25), isls-templates (C26), isls-foundry (C27)
  external:
    glyph-ir, glyph-embed, glyph-canon, glyph-q16
    (path deps to babylon-compiler workspace)

isls-harness (C10) -- MODIFIED
  new deps: isls-forge, isls-oracle, isls-templates,
            isls-foundry, isls-gateway
  new bench functions: B16-B24
\end{lstlisting}


\section{Test Count}

\begin{tabular}{l r r}
\toprule
\textbf{Phase} & \textbf{New Tests} & \textbf{Cumulative} \\
\midrule
Phase 1--9 & 326 & 326 \\
Phase 10A (C28 Babylon Bridge) & 12 & 338 \\
Phase 10B (Benchmarks B16--B24) & 9 & $\ge$ 347 \\
\bottomrule
\end{tabular}


\section{Completion Criteria}

\begin{enumerate}[nosep]
  \item ArtifactIR converts to glyph-ir deterministically.
  \item H5 embedding validates structural coherence before codegen.
  \item Six codegen backends produce scaffolding from IR structure.
  \item Oracle fills function bodies; backends assemble final files.
  \item Cross-language type consistency from single IR node.
  \item OpsBackend and DocsBackend produce complete files without
        Oracle.
  \item Six full-stack templates (T11--T16) produce multi-language
        projects.
  \item Babylon crates integrated as workspace dependencies.
  \item B16--B24 benchmarks implemented and running.
  \item HTML report includes Section~11 for generative benchmarks.
  \item \texttt{isls bench --suite generative} works.
  \item Total tests $\ge 347$, zero warnings.
  \item \textbf{Backends write scaffolding. Oracle writes logic.
        Benchmarks measure everything.}
\end{enumerate}


\vfill
\begin{center}
\rule{0.5\textwidth}{0.4pt}\\[0.5cm]
{\small Generated for \ISLS{} Extension Phase 10 v1.0.0 (Revised v2).\\[0.2cm]
Honest inventory. Honest architecture. Honest benchmarks.\\
Babylon validates structure. Backends write scaffolding.\\
Oracle writes logic. Foundry verifies. Benchmarks measure.\\
Nothing is claimed that does not exist.}
\end{center}

\end{document}
